\providecommand{\devnum}[1]{#1}

\section{Traces and Workload Characterisation}
\label{sec:data}

\subsection{The Traces}

The validation scenario is an educational one: two programming-assignment judging
platforms supply the demand traces, and four further workloads bound the claim from
outside that setting. What each trace can be used for is decided by whether its arrival
times are observed or reconstructed, whether the queue wait is recorded, and how long a job may run there
(Table~\ref{tab:traces}).

% Wider than the MDPI text block, so it is set across the full page with the
% template's adjustwidth/\extralength mechanism; \fulllength is that measure.
\begin{table}[htbp]
\begin{adjustwidth}{-\extralength}{0cm}
\centering
\caption{The traces, their role and what they record. ``Arrivals'' says whether the arrival
instants are recorded, reconstructed proxies, or drawn; ``Waits'' whether the queue wait of each job is
itself recorded.}
\label{tab:traces}
\small
\begin{tabular}{>{\raggedright\arraybackslash}p{0.13\fulllength}%
                >{\raggedright\arraybackslash}p{0.17\fulllength}%
                >{\raggedright\arraybackslash}p{0.20\fulllength}%
                >{\raggedright\arraybackslash}p{0.16\fulllength}cc}
\hline
Trace & Role & Size & Time fields & Arrivals & Waits \\
\hline
CodeBench~\cite{codebench_dataset} &
main demand trace &
\devnum{4{,}611{,}325} blocks, \devnum{863{,}149} graded submissions in the development pool &
one stamp per block, the end of the run &
start proxy & no \\
ACcoding~\cite{accoding2024} &
second judging platform &
\devnum{4{,}046{,}652} rows, \devnum{3{,}237{,}322} non-sealed &
none; submission id only &
no, drawn & no \\
Serverless invocations \cite{azure2021trace} &
method applies outside education &
\devnum{1{,}980{,}951} invocations over \devnum{14.000} d &
modified end stamp and duration &
start proxy & no \\
CI pool A \cite{taskcluster_docs,treeherder_docs} &
simulator validation, light tail &
\devnum{191{,}970} runs over \devnum{21} d &
created, scheduled, started, resolved &
yes & yes \\
CI pool B &
second validation pool &
\devnum{30{,}702} runs over \devnum{6.99} d &
as above & yes & yes \\
Compute farm \cite{netbatch2012trace,shai2014netbatch} &
non-applicable case &
\devnum{9{,}054{,}066} jobs over \devnum{30.00} d, \devnum{8{,}441{,}207} serial with positive run time &
submit, wait, run &
yes & yes \\
OULAD~\cite{oulad2017} &
calendar-driven load elsewhere &
\devnum{9} course-presentations, daily counts &
day &
--- & --- \\
\hline
\end{tabular}
\end{adjustwidth}
\end{table}
% Sources of Table~\ref{tab:traces}:
%   CodeBench sizes: evidence/codebench_audit/out_tail_and_counts.txt (4,611,325 events;
%     1,398,801 submits; 3,212,524 tests) and evidence/accoding_v2/out_accoding_v2.txt
%     stage compare section B (861,478 graded submissions with C_cap > 0; 863,149 is the
%     block count before the C == 0 drop, evidence/codebench_service_v2/out_service_v2.txt).
%   ACcoding sizes: evidence/accoding_v2/out_accoding_v2.txt, stages selftest and audit 1b.
%   Serverless: evidence/cross_domain/out_audit_azure.txt (rows, span).
%   CI pools: evidence/firefox_ci/out_SUMMARY.txt section A and out_simval.txt (spans).
%   Compute farm: evidence/cross_domain/out_audit_netbatch.txt.
%   OULAD: evidence/oulad/out_load.txt (header line, nine 2013 course-presentations).

Both CI pools and the compute farm carry a recorded queue wait. We validate the simulator
against the CI pools, whose capacity is a countable set of physical machines that changes
only on reimaging or quarantine. The second judging platform has no timestamp
column, only a submission id, so job order inside a group and job sizes are real while
every arrival instant is drawn. The correlation between a group's first id and its
recorded start time is \devnum{0.9995} over \devnum{469} groups, but in \devnum{207} of
\devnum{468} consecutive transitions the id ranges overlap, so the id is an ordering and
not a clock. Its results are a check that the ordering mechanism transfers, not a wait in
seconds.

\subsection{What the Fields of the Main Trace Mean}

Several fields of that trace had to be interpreted before a queueing model could use
them. Four readings carry the rest of the paper; Supplementary Section~S2 gives the
audit that establishes them, together with the release-statistics reconciliation and the
two truncated archives.

The timestamp on an execution block is the moment the result was written, not the moment
the code was received, so we take $\mathrm{done}_j$ to be the recorded stamp and
reconstruct the arrival as $a_i = t_i - C_i$, carrying the opposite reading through the
pipeline as a sensitivity with its own features and models. Stamps are recorded to the
whole second while costs are recorded to the microsecond, which opens a leak that we
close by jittering each stamp deterministically: a classifier given only the fractional
part of the gap to the previous record separates heavy from light jobs with AUROC
\devnum{0.8091} on the whole-second clock and \devnum{0.4499} after the jitter, while
the forward models of the two clocks stay within \devnum{0.008} of AUROC of each other.
The observed regime has an upper ceiling of \devnum{60} s; this is not evidence of
a documented enforced timeout. The last run to reach a longer time ends at
\devnum{2019-08-28}; the periods before that date stay out of the primary pool, and
where used at all they carry the later limit, $C^{\mathrm{cap}} = \min(C, 60)$
(Table~\ref{tab:tails}). Groups of runs by the same user that end in the same second are
\devnum{7.53\%} of blocks and \devnum{26.28\%} of capped work; they overlapped and
stretched each other, so we neither delete them silently nor treat each as an
independent arrival, and the scheduling trace is rebuilt twice, once with each group
collapsed into one job and once with the groups removed (Table~\ref{tab:sens}).

\subsection{What the Source Platform Identifies}

The platform that produced the main trace gives each logged-in student a dedicated
container and runs their code inside it. Its logs do not identify the waiting times
of a shared judging queue. We use recorded durations and reconstructed start-time
proxies as inputs to a hypothetical shared pool; its capacity, queue and service
order are our design. Neither this construction nor the absence of shared-queue
fields establishes that all source-platform work experienced no waiting.

Three things support the design. The first is that a fixed pool of judging hosts on one
queue is a standard deployment. A widely used contest judging
system recommends ``one judgehost per 20 teams''~\cite{domjudge_manual} on dedicated
machines, and the autograder study of Section~\ref{sec:intro} runs a configurable number
of concurrent graders, replays the busiest hours before a shared deadline and measures the
queue wait~\cite{peveler2019sandbox}. That replay compresses its window fourfold and
sixteenfold; we raise load by overlaying whole class-terms instead
(Section~\ref{sec:overlay}), for the reason Section~\ref{sec:limits} gives. Dedicated
per-entity capacity is also known to sit
idle in the workloads this one resembles. In the serverless characterisation behind our
invocation trace, \devnum{45\%} of applications are invoked once an hour or less and
\devnum{81\%} once a minute or less, so the authors report that the cost of keeping them
warm relative to their billable execution time can be prohibitively
high~\cite{shahrad2020wild}; and a cluster manager that measured the alternative directly
reports that splitting each large user's workload into its own cell would need
\devnum{2} to \devnum{16} times as many cells and \devnum{20} to \devnum{150\%} more
machines~\cite{verma2015borg}. Neither source is about judging, and neither licenses a
number of our own.

The second is a measurement on the same demand. Replaying the overlay of
Section~\ref{sec:overlay} under the dedicated design, the 99th percentile of containers
alive at once is \devnum{1{,}636} and each container executes code for \devnum{0.12\%} of
the time it is alive; since only \devnum{11\%} to \devnum{35\%} of logins have a logout
written after them, session boundaries are reconstructed, and across six reconstructions
that percentile lies between \devnum{1{,}003} and \devnum{2{,}525} containers and the
utilisation between \devnum{0.07\%} and \devnum{0.17\%}. A shared pool holds the 99th
percentile of deadline-window waits under \devnum{1} s with \devnum{14} servers under
FCFS and \devnum{13} under the predictor-ordered policy. Pooling is worth two orders of
magnitude here and the scheduler is worth one server (\devnum{14} to \devnum{13} at
\devnum{1} s, \devnum{11} to \devnum{10} at \devnum{5} s, \devnum{8} to \devnum{7} at
\devnum{30} s). It does not support a claim that the platform ever queued, the pooled
side being simulated, nor that the containers are over-provisioned, since they also run
the editing session the pool does not replace, nor that any of it is a cost saving.

The third is that the simulator reproduces waits somebody else recorded. The CI logs carry
the queue wait of every run, so we replay the recorded arrivals and services through our
own $k$-server non-preemptive kernel. Two quantities are fitted: the effective capacity
$k_{\mathrm{eff}}$, the shortfall against the listed machine count being machines that are
up but not claiming work, and a setup term, the median short gap between consecutive runs
on the same machine, added to every service because a machine is not reusable the instant
a run resolves. Fitting both on the window and then reporting agreement on that same
window would be circular, so we refit both on one chronological part of each window and
test on the other, five splits over the two pools.
This split is internal to the recorded window and is held out of the fit only; it
is not sealed in the sense of Section~\ref{sec:sealed}, and the whole window has
been read.

The capacity is stable across splits. On pool~A the three splits choose \devnum{155},
\devnum{155} and \devnum{157} of \devnum{173} listed machines, and on pool~B the two
splits choose \devnum{96} and \devnum{92} of \devnum{105}; the setup term moves between
\devnum{54.2} and \devnum{55.3} s on the first and between \devnum{126.2} and
\devnum{128.6} s on the second. The objective is sharp in $k$. With those parameters
frozen, the held-out part comes
back at \devnum{0.90} to \devnum{1.14} times the recorded mean wait, \devnum{0.86} to
\devnum{1.16} at the 99th percentile and \devnum{0.89} to \devnum{1.28} at the 99.9th,
with a per-hour mean-wait correlation of \devnum{0.936} to \devnum{0.981}, a per-job rank
correlation of \devnum{0.940} to \devnum{0.963}, and a median per-job absolute error of
\devnum{16} to \devnum{84} s.

The 90th percentile does not survive the held-out test. It lands at \devnum{0.81} to
\devnum{1.21}
of the recorded one, and on the larger pool it is below the recorded value on all three
splits while the 99th is above it, so the simulated distribution is too steep in its upper
middle; the two-term objective hides this in sample by trading the mean against the p90,
so we fit out of sample and say so. What the fit absorbs is nonetheless a
property of the pool and no artefact of fitting: told instead that every listed
machine is available, the same simulator underestimates the held-out mean wait by a factor
of \devnum{1.3} to \devnum{2.6} and its per-hour correlation falls to \devnum{0.75}.
Supplementary Section~S2 gives the step-of-one sharpness check and reads the circularity
of the original whole-window fit off the two pools directly.
% Sources of the three paragraphs above:
%   evidence/firefox_ci_holdout/out_SUMMARY.txt, section A (k_eff step-1 per split and the
%   setup term on the fit part: osx 155/155/157 of 173, talos 96/92 of 105; setup 54.2-55.3 s
%   and 126.2-128.6 s), section B (the HELD OUT rows: mean x, p99 x, p99.9 x, rho_h, sp,
%   med|e|, and the "no fit k = nominal machines" rows for the 1.3-2.6x and the 0.75),
%   section E (5-8 values of k inside 10% of the recorded mean) and section F (the p90
%   range 0.81-1.21 and the circularity comparison 1.14 vs 0.99 on the talos reverse split).
%   The same rows are in holdout.csv.  Pool A = releng-hardware/gecko-t-osx-1500-m4
%   (173 machines, 21 d), pool B = releng-hardware/gecko-t-linux-talos-2404 (105, 7 d),
%   matching Table~\ref{tab:traces}.  The sharpness check and the circularity read-off
%   moved to supplementary.tex, Supplementary Section~S2 (referee item G9).

\paragraph{Capacity and physical checks.}
The supplementary capacity sweep retains the complete offered stream and varies
the number of servers; it uses the archived embedded Tweedie fit without refitting.
The physical experiment executes FCFS, SPJF-E and a guard on the complete selected
development window, with timed slot occupancy and recorded enqueue, dispatch and
completion events. It tests an executing rule but does not transport submitted
programs to new hardware. The selected empty-start window is distinct from a
full-trace replay; its guarded p99 is worse and its job-level ideal-simulator
agreement fails (Section~\ref{sec:exp_physical}, Supplementary Section~S6).

\subsection{Building the Shared Pool}
\label{sec:overlay}

A single class-term does not load one server, so we overlay several. The pool takes only the development
terms in the observed \devnum{60} s regime, \devnum{40} class-terms in all, and
replicates them \devnum{44} times, each copy aligned by weekday and hour and shifted by a
whole number of weeks drawn per copy. Overlaying whole terms rather than compressing
inter-arrival times preserves the daily and weekly cycle and each job's position relative
to its own deadline. The result is
\devnum{17{,}634{,}760} jobs over \devnum{30} weeks and \devnum{208.8} days. The three
load levels share that trace and differ only in the pool size. Four of the five overlays
take \devnum{$k = 8/5/4$} and the fifth takes \devnum{$7/5/4$}, its work per busy hour
being the lowest of the five; on the first overlay that gives busy-hour utilisations of
\devnum{0.529}, \devnum{0.847} and \devnum{1.059}, and \devnum{0.502}, \devnum{0.783} and
\devnum{0.979} averaged over the five. A
deadline window is the \devnum{24} hours before an assessment's own end, and
\devnum{4{,}219{,}556} jobs arrive inside one. Five overlays are week-shifts of the same
terms, so they are one construction and not five platforms, and our intervals resample
whole weeks, which covers time and leaves the choice of terms uncovered. A trace built
from the last development term alone is
not available this way, since it would need \devnum{190} copies of its \devnum{8}
class-terms to load a pool of four.

\subsection{What the Load Looks Like}

Arrivals follow the course calendar. Under grouped-class cross-validation,
predicting executions per class \devnum{168} hours ahead gives log-scale RMSE
\devnum{1.243} with lag features alone, \devnum{1.159} with calendar features
alone, and \devnum{1.112} with both. A paired bootstrap over class-week blocks
gives a combined-minus-lag-only difference of
\devnum{$-0.1307$ $[-0.1395, -0.1213]$}.
% Source: evidence/codebench/out_codebench.txt, P4 grouped-class-CV h=168
% (lgbm_lags, lgbm_cal_only, lgbm_lags_cal) and paired CV h=168 rows.
 The corroborating trace shows the same shape
elsewhere: clicks in the three days up to an assessment due date are \devnum{1.29} to
\devnum{2.95} times the other test days across its nine course-presentations, and a
seven-day-ahead peak warning on distinct active students goes from F1 \devnum{0.000} under
a moving-average baseline to \devnum{0.593} once calendar features are added.

Costs are concentrated in a small fraction of jobs on every workload except the CI
pools, which motivates the empirical cost-concentration screen (Table~\ref{tab:tails}).

\begin{table}[htbp]
\centering
\caption{Share of executed work carried by the longest jobs, and provenance of the
analysis bound $L$. An analysis cap is applied to the replay, not evidence of a
platform-enforced timeout.}
\label{tab:tails}
\small
\begin{tabular}{lccp{0.30\textwidth}}
\hline
Trace & Top \devnum{1\%} & Top \devnum{5\%} & Bound and provenance \\
\hline
CodeBench, \devnum{60} s regime & \devnum{0.375} & \devnum{0.432} & \devnum{60} s, observed regime ceiling; analysis cap \\
CodeBench, no-limit periods, raw & \devnum{0.941} & \devnum{0.960} & none \\
CodeBench, no-limit periods, capped & \devnum{0.494} & \devnum{0.658} & \devnum{60} s, imposed analysis truncation \\
ACcoding, measured cost & \devnum{0.275} & \devnum{0.628} & \devnum{30.0} s, observed maximum of measured test-cost sum \\
ACcoding, with the assumed fixed term & \devnum{0.122} & \devnum{0.290} & \devnum{30.5} s, observed maximum after assumed overhead \\
Serverless invocations & \devnum{0.443} & \devnum{0.879} & \devnum{286.8} s, training p99.9 analysis cap \\
Compute farm & \devnum{0.372} & \devnum{0.628} & \devnum{493{,}857} s, training p99.9 analysis cap \\
CI pool A & \devnum{0.034} & \devnum{0.107} & \devnum{5{,}066} s, training p99.9 of run plus setup; clipped \\
CI pool B & \devnum{0.052} & \devnum{0.159} & \devnum{4{,}140} s, training p99.9 of run plus setup; clipped \\
\hline
\end{tabular}
\end{table}
% Sources of Table~\ref{tab:tails}:
%   CodeBench rows: evidence/codebench_audit/out_tail_and_counts.txt (eras 2020-ERE..2022-2
%     and 2018-1..2019-2, columns top1_share / top5_share, raw C and C_cap = min(C,60)).
%   ACcoding rows and both per-job limits: evidence/accoding_v2/out_accoding_v2.txt,
%     stage compare section B (time_cost and 0.5 + time_cost columns).
%   Serverless and compute farm: evidence/cross_domain/out_cross_domain_table.txt section A
%     (L used column) and out_audit_azure.txt / out_audit_netbatch.txt (work-share lines).
%   CI pools: evidence/firefox_ci/out_SUMMARY.txt section A ("share of total work carried
%     by the top" and "L' (train p99.9 of service)").

\paragraph{Arrival and service provenance across traces.}
CodeBench's jittered result stamp minus the recorded duration is an execution-start
proxy, not a measured receipt time. ACcoding has no submission clock: arrivals are
drawn, including the Poisson comparison and the deadline-shaped construction; its
submission ids determine ordering, not elapsed time. Its measured cost sums test
case execution times, each subject to its problem limit, and the replay adds an
assumed \devnum{0.5} s overhead. Thus neither the observed \devnum{30.0} s sum nor
the resulting \devnum{30.5} s maximum is a verified whole-job timeout.
% Source: evidence/accoding_v2/out_accoding_v2.txt, field audit and compare section B.

For Azure, the replay arrival is the documented end timestamp minus duration,
again a start proxy; the publisher explicitly states that invocation timestamps
were modified from the production trace \cite{azure2021trace}. The Netbatch replay
uses the recorded SWF submission time and run duration. Their bounds in
Table~\ref{tab:tails} are training p99.9 truncations, not observed platform timeouts.
Both FCFS and every reordered policy receive the same truncated service sequence.
The truncation affects \devnum{1{,}187} Azure jobs and \devnum{7{,}696} Netbatch jobs
in the analysed pools; it removes \devnum{3.18\%} and \devnum{2.16\%} of raw work.
% Sources: evidence/cross_domain/meta_azure.csv and meta_netbatch.csv;
% evidence/cross_domain/out_leaktest_azure.txt and out_leaktest_netbatch.txt,
% clipped counts; out_audit_azure.txt and out_audit_netbatch.txt, work removed.

In the CI pools, arrival is the recorded scheduled instant, when dependencies have
cleared and a worker may claim the run; recorded wait is started minus scheduled.
Service is resolved minus started plus a fitted same-worker setup term. The analysis
first adds setup and then clips this occupancy at the training p99.9 $L$;
longer jobs are retained with truncated service, not excluded. This clips
\devnum{227} runs in pool~A and \devnum{27} in pool~B. FCFS and all comparison
policies use exactly that same capped occupancy sequence. The source scheduler's
priority-then-FIFO rule and operational capacity are distinct from the counterfactual
FCFS model used for the guarantee.
% Sources: evidence/firefox_ci/out_leaktest_releng-hardware__gecko-t-osx-1500-m4.txt and
% out_leaktest_releng-hardware__gecko-t-linux-talos-2404.txt;
% evidence/firefox_ci/meta_releng-hardware__gecko-t-osx-1500-m4.csv and
% meta_releng-hardware__gecko-t-linux-talos-2404.csv; fci_build.py and fci_sim.py.

An enforced $L$ in a deployed service would instead kill an unfinished job at the
limit and charge that executed work to the same-input FCFS comparison. The killed
job would not yield its intended result; these replays do not estimate the resulting
completion-quality cost or users' resubmission response. The table therefore
supports finite-input bounds under its stated transformations, not a future
operational timeout claim.

\subsection{Sealed Data}
\label{sec:sealed}

Three of the traces carry a part that was sealed before the method was designed.
On the main trace it is three terms, \devnum{2023-1}, \devnum{2023-2} and
\devnum{2024-1}; on the second judging platform it is the last \devnum{20\%} of
the submission-id range, \devnum{809{,}330} rows, of which only the row count and
the first id have been read; on the corroborating trace it is the \devnum{2014}
academic year. Only the main trace's sealed terms will be opened for this paper, once,
after the method, its parameters and the primary metric have been frozen; the frozen
pipeline covers those three terms and nothing else, and the sealed parts of the
other two traces stay shut and are reported on by neither. The serverless trace,
the compute farm and the two continuous-integration pools have no sealed part,
so every figure we give for them is development data. No cost distribution,
prediction or scheduling result has been computed on any sealed part, and every
access to one is entered in a ledger in the repository. Every number in this
section and in Sections~\ref{sec:prediction} and~\ref{sec:scheduling} is a
development number. The sealed evaluation remains pending in this manuscript;
Section~\ref{sec:exp_prov} states the current output scope.
% Source: prechecks/heldout_scope/heldout_scope.md, G1--G6; evidence/accoding_v2/
% out_accoding_v2.txt, audit/selftest; evidence/codebench_service_v2/out_service_v2.txt,
% sealed-term exclusions; evidence/oulad/out_load.txt, development presentations.
